\documentclass[aspectratio=169,11pt]{beamer}

% Compatibilidad con compilaciones mediante pdfLaTeX.
\usepackage[utf8]{inputenc}
\usepackage[T1]{fontenc}
\usepackage[spanish]{babel}

% Matemáticas y diagramas.
\usepackage{amsmath,amssymb,mathtools}
\usepackage{tikz}
\usetikzlibrary{matrix,positioning}
\usepackage{tikz-cd}

% Tema docente, incorporado como submódulo en beamerteaching/.
\makeatletter
\def\input@path{{beamerteaching/}}
\makeatother
\usetheme{Teaching}
\titlegraphic{%
  \includegraphics[width=0.20\paperwidth]{beamerteaching/assets/logo-us.png}%
}

\title{Álgebra Lineal y Geometría I}
\subtitle{Grado en Matemáticas}
\author{Fernando Muro}
\institute{Universidad de Sevilla}
\date{Curso 2026--2027}

% Comandos generales de la asignatura.
\newcommand{\R}{\mathbb{R}}
\newcommand{\C}{\mathbb{C}}
\newcommand{\K}{\mathbb{K}}
\newcommand{\Mat}{\operatorname{Mat}}
\newcommand{\rg}{\operatorname{rg}}
\newcommand{\im}{\operatorname{im}}
\newcommand{\id}{\operatorname{id}}

\begin{document}

\begin{frame}
  \titlepage
\end{frame}

\section{Sistemas de ecuaciones lineales, matrices y determinantes}

\begin{frame}{Cuerpos}

Trabajaremos con matrices cuyos coeficientes pertenecen a un \alert{cuerpo} \(K\).

\medskip

Para nosotros, un cuerpo es un conjunto de ``números'' en el que podemos:

\begin{itemize}
    \item sumar,
    \item restar,
    \item multiplicar,
    \item dividir por cualquier número distinto de cero.
\end{itemize}

\medskip

Ejemplos de cuerpos:
\[
\mathbb{Q},
\qquad
\mathbb{R},
\qquad
\mathbb{C}.
\]

\medskip

En cambio, \(\mathbb{N}\) y \(\mathbb{Z}\) no son cuerpos:

\begin{itemize}
    \item en \(\mathbb{N}\) no siempre podemos restar:
    \[
    2-3\notin\mathbb{N};
    \]
    \item en \(\mathbb{Z}\) no siempre podemos dividir por un número no nulo:
    \[
    \frac{1}{2}\notin\mathbb{Z}.
    \]
\end{itemize}

\end{frame}
% \input{temas/01-sistemas-matrices-determinantes/ejercicios/...}

\input{temas/02-espacios-vectoriales/tema02}
\input{temas/03-subespacios/tema03}
\input{temas/04-homomorfismos/tema04}
\input{temas/05-forma-canonica/tema05}
\input{temas/06-espacio-vectorial-euclideo/tema06}
\input{temas/07-teoremas-espectrales/tema07}
\input{temas/08-espacios-afines/tema08}
\input{temas/09-aplicaciones-afines/tema09}
\input{temas/10-espacios-afines-euclideos/tema10}
\input{temas/11-movimientos/tema11}

\end{document}
